\documentclass[12pt,a4paper,titlepage]{article}

\usepackage{xcolor}
\usepackage{minted}
\usepackage{graphicx}
\usepackage{titlesec}
\usepackage[portuguese]{babel}
\usepackage[colorlinks=true,allcolors=teal]{hyperref}

\graphicspath{{./img/}}
\newcommand{\sectionbreak}{\clearpage}

\title{
Arquitetura de Computadores II
\\Medição e Análise de Desempenho
\\[1.5em]
\centering
\includegraphics[scale=0.8]{puc.png}
}
\author{
    Fernando de Lima Gonçalves\\
    João Pedro Mesquita Salgado Gil\\
    Laura Martins de Souza\\
    Samuel Cândido Danelon da Cruz
}
\date{Belo Horizonte\\10/2026}

\begin{document}
\maketitle
\tableofcontents
\clearpage

% Conteúdo
\section{Introdução}
\subsection{Objetivo Geral}
\textit{
    "Compreender na prática os desafios e as ferramentas utilizadas para medir
    o tempo de execução e analisar o desempenho arquitetural de sistemas computacionais,
    fugindo das armadilhas de métricas teóricas isoladas."
}

\subsection{Atividades}
As atividades a serem desenvolvidas foram particionadas em 3 grupos/partes:
\begin{enumerate}
    \item 
        \textbf{Metodologia de Medição de Tempo de Execução}
        \begin{enumerate}
            \item Pesquisa e Instrumentação
            \item Execução e Coleta de Dados
            \item Ambiente de Execução
            \item Análise Crítica
        \end{enumerate}
    \item \textbf{Análise de Desempenho em Nível Arquitetural com “perf”}
    \item \textbf{Avaliação Comparativa com Benchmarks Padrão (SPEC)}
\end{enumerate}

Para realização das atividades utilizou-se dos algoritmos nas linguagens:
\textbf{\textit{C, Python e Java}}.

\section{Desenvolvimento}
O Algoritmo designado ao grupo foi: \textit{Somatório de elementos de um vetor.}

\subsection{Algoritmos em C}
\subsubsection{clock()}
\inputminted[fontsize=\small,linenos,frame=single,breaklines=true]{c}{../C/clock.c}

\subsubsection{time()}
\inputminted[fontsize=\small,linenos,frame=single,breaklines=true]{c}{../C/time.c}

\subsubsection{clock\_gettime()}
\inputminted[fontsize=\small,linenos,frame=single,breaklines=true]{c}{../C/gettime.c}

\subsection{Pesquisa e Instrumentação}
\begin{itemize}
    \item \textit{clock():} Mede o tempo de CPU consumido pelo processo. Útil para avaliar processamento, sem considerar diretamente períodos de espera.
    \item \textit{time():} Obtém o horário atual, normalmente em segundos. Possui baixa resolução, sendo inadequado para medir algoritmos rápidos.
    \item \textit{clock\_gettime():} Obtém o tempo de um relógio especificado, com resolução elevada. 
        Com \textit{CLOCK\_MONOTONIC}, permite medir intervalos curtos sem sofrer ajustes do relógio do sistema.
\end{itemize}

\subsection{Medições}
O grupo optou pela utilização de diferentes computadores mas garantindo que o 
interpretador/máquina virtual sejam iguais ou mais próximos possíveis.

\begin{center}
    \begin{tabular}{ || c | c | c | c || }
        \hline
        \multicolumn{4}{||c||}{\textbf{Computadores}}\\
        \hline \hline
        \textit{Processador} & \textit{RAM} & \textit{SO} & \textit{Compilador} \\ [0.5ex]
        \hline
        Apple M1 & 8gb & Tahoe 26.6.2 & GCC 16.2 \\
        \hline
        AMD Ryzen 9 5900x & 16gb & Linux Fedora 7.2.11 & GCC 16.2 \\
        \hline
    \end{tabular}
\end{center}

\clearpage\subsubsection{C: clock()}
\begin{center}
    \begin{tabular}{||c|c|c|c||}
        \hline
        \multicolumn{4}{||c||}{\textbf{C - clock(): macOS}}\\
        \hline \hline
        Medições & \textit{Instância: 100} & \textit{Instância: 1000} & \textit{Instância: 10000}\\
        \hline
        1ª medição & 0.000006[s] & 0.000010[s] & 0.000027[s]\\
        \hline
        2ª medição & 0.000008[s] & 0.000010[s] & 0.000025[s]\\
        \hline
        3ª medição & 0.000006[s] & 0.000013[s] & 0.000026[s]\\
        \hline
        4ª medição & 0.000011[s] & 0.000009[s] & 0.000030[s]\\
        \hline
        5ª medição & 0.000004[s] & 0.000011[s] & 0.000045[s]\\
        \hline
        6ª medição & 0.000011[s] & 0.000008[s] & 0.000030[s]\\
        \hline
        7ª medição & 0.000009[s] & 0.000010[s] & 0.000029[s]\\
        \hline
        8ª medição & 0.000009[s] & 0.000007[s] & 0.000027[s]\\
        \hline
        9ª medição & 0.000008[s] & 0.000008[s] & 0.000035[s]\\
        \hline
        10ª medição & 0.000007[s] & 0.000006[s] & 0.000033[s]\\
        \hline
        \textbf{Média} & \textbf{0.0000079[s]} & \textbf{0.0000092[s]} & \textbf{0.0000307[s]}\\
        \hline
    \end{tabular}
\end{center}

\subsubsection{C: time()}
\begin{center}
    \begin{tabular}{||c|c|c|c||}
        \hline
        \multicolumn{4}{||c||}{\textbf{C - time(): macOS}}\\
        \hline \hline
        Medições & \textit{Instância: 100} & \textit{Instância: 1000} & \textit{Instância: 10000}\\
        \hline
        1ª medição & 0.000000[s] & 0.000000[s] & 0.000000[s]\\
        \hline
        2ª medição & 0.000000[s] & 0.000000[s] & 0.000000[s]\\
        \hline
        3ª medição & 0.000000[s] & 0.000000[s] & 0.000000[s]\\
        \hline
        4ª medição & 0.000000[s] & 0.000000[s] & 0.000000[s]\\
        \hline
        5ª medição & 0.000000[s] & 0.000000[s] & 0.000000[s]\\
        \hline
        6ª medição & 0.000000[s] & 0.000000[s] & 0.000000[s]\\
        \hline
        7ª medição & 0.000000[s] & 0.000000[s] & 0.000000[s]\\
        \hline
        8ª medição & 0.000000[s] & 0.000000[s] & 0.000000[s]\\
        \hline
        9ª medição & 0.000000[s] & 0.000000[s] & 0.000000[s]\\
        \hline
        10ª medição & 0.000000[s] & 0.000000[s] & 0.000000[s]\\
        \hline
        \textbf{Média} & \textbf{0.000000[s]} & \textbf{0.000000[s]} & \textbf{0.000000[s]}\\
        \hline
    \end{tabular}
\end{center}

\subsubsection{C: clock\_gettime()}
\begin{center}
    \begin{tabular}{||c|c|c|c||}
        \hline
        \multicolumn{4}{||c||}{\textbf{C - time(): macOS}}\\
        \hline \hline
        Medições & \textit{Instância: 100} & \textit{Instância: 1000} & \textit{Instância: 10000}\\
        \hline
        1ª medição & 0.000000[s] & 0.000002[s] & 0.000030[s]\\
        \hline
        2ª medição & 0.000001[s] & 0.000003[s] & 0.000024[s]\\
        \hline
        3ª medição & 0.000000[s] & 0.000003[s] & 0.000037[s]\\
        \hline
        4ª medição & 0.000001[s] & 0.000002[s] & 0.000031[s]\\
        \hline
        5ª medição & 0.000000[s] & 0.000003[s] & 0.000024[s]\\
        \hline
        6ª medição & 0.000000[s] & 0.000003[s] & 0.000031[s]\\
        \hline
        7ª medição & 0.000001[s] & 0.000003[s] & 0.000031[s]\\
        \hline
        8ª medição & 0.000001[s] & 0.000002[s] & 0.000024[s]\\
        \hline
        9ª medição & 0.000001[s] & 0.000003[s] & 0.000031[s]\\
        \hline
        10ª medição & 0.000005[s] & 0.000003[s] & 0.000031[s]\\
        \hline
        \textbf{Média} & \textbf{0.0000010[s]} & \textbf{0.0000027[s]} & \textbf{0.0000294[s]}\\
        \hline
    \end{tabular}
\end{center}

\section{Conclusão}
Com base nos testes desenvolvidos foi possível perceber pontos positivos e negativos dos métodos de avaliação de desempenho.
\begin{itemize}
    \item \textit{time():} Simples implementação mas fortemente limitado devido a intervalo definido em segundos.
    \item \textit{clock():} Medição eficiente para tempo consumido pela CPU durante processo avaliado porém sofre influência da carga de processamento.
    \item 
        \textit{clock\_gettime():} Permite medições de alta resolução e com inclusão CLOCK\_MONOTONIC não sofre interferência do relógio do sistema.
        Seu ponto negativo pode se resumir a sua dependência do sistema operacional e hardware.
\end{itemize}
\end{document}
